% =====================================================================
% config.tex - Cấu hình toàn cục (margin, header, footer, colors, styles)
% =====================================================================

% ======================== CẤUHÀNH PAGE =============================
\onehalfspacing                      % Giãn dòng 1.5

% ======================== HEADER & FOOTER =============================
\pagestyle{fancy}
\fancyhf{}                           % Xóa mặc định
\fancyhead[L]{\nouppercase{\leftmark}}    % Chương hiện tại bên trái
\fancyhead[R]{\thepage}                   % Số trang bên phải
\renewcommand{\headrulewidth}{0.4pt}      % Đường kẻ header

% ======================== ĐỊNH NGHĨA MÀU SẮC =============================
\definecolor{codegreen}{rgb}{0,0.6,0}
\definecolor{codegray}{rgb}{0.5,0.5,0.5}
\definecolor{codeblue}{rgb}{0,0,0.8}
\definecolor{backcolour}{rgb}{0.95,0.95,0.92}

% ======================== CẤU HÌNH LISTINGS (CODE) =============================
\lstdefinestyle{sqlstyle}{
    backgroundcolor=\color{backcolour},
    commentstyle=\color{codegreen},
    keywordstyle=\color{codeblue}\bfseries,
    numberstyle=\tiny\color{codegray},
    stringstyle=\color{codegreen},
    basicstyle=\ttfamily\footnotesize,
    breakatwhitespace=false,
    breaklines=true,
    captionpos=b,
    keepspaces=true,
    numbers=left,
    numbersep=5pt,
    showspaces=false,
    showstringspaces=false,
    showtabs=false,
    tabsize=2,
    language=SQL,
    frame=single,
    rulecolor=\color{black}
}
\lstset{style=sqlstyle}

% ======================== CẤU HÌNH TIÊU ĐỀ =============================
\setcounter{secnumdepth}{3}          % Đánh số đến subsubsection
\setcounter{tocdepth}{3}             % Hiển thị đến subsubsection trong mục lục

% ======================== LIÊN KẾT HYPERREF =============================
\hypersetup{
    colorlinks=true,
    linkcolor=blue,
    filecolor=magenta,
    urlcolor=cyan,
    pdftitle={Thiết Kế Cơ Sở Dữ Liệu},
    pdfauthor={Nhóm Thiết Kế},
    bookmarksnumbered=true
}

% ======================== LỆNH TÙY CHỈNH =============================

% Lệnh cho từ điển dữ liệu
\newcommand{\datadict}[4]{%
    \textbf{#1} & #2 & #3 & #4 \\ \hline
}

% Lệnh cho bảng thuật ngữ
\newcommand{\term}[2]{%
    \textbf{#1} & #2 \\ \hline
}

% Tên entity/table có thể chứa dấu gạch dưới mà không cần escape thủ công.
% \detokenize ngăn ký tự "_" làm LaTeX chuyển nhầm sang chế độ toán.
\newcommand{\entity}[1]{\texttt{\detokenize{#1}}}

% Lệnh cho các ký hiệu quan trọng
\newcommand{\highlight}[1]{\textbf{\textcolor{codeblue}{#1}}}

% ======================== CẤU HÌNH PARAGRAPH =============================
\parindent=1cm                       % Thụt đầu dòng
\renewcommand{\baselinestretch}{1.5} % Giãn dòng 1.5
